\documentclass[tikz,border=6pt]{standalone}
\usepackage{ctex}
\usepackage{amsmath}
\usetikzlibrary{calc, arrows.meta}

\begin{document}
\begin{tikzpicture}[scale=1.0, thick]

% 含参直线扫描圆锥曲线示意图
% 椭圆 x^2/4 + y^2 = 1，一组斜率不同的直线过点 P(0, 0.5)

\def\a{2.0}
\def\b{1.0}

% 坐标轴
\draw[->, gray, thin] (-2.8, 0) -- (2.8, 0) node[right, font=\footnotesize] {$x$};
\draw[->, gray, thin] (0, -1.5) -- (0, 2.0) node[above, font=\footnotesize] {$y$};
\node[font=\footnotesize, below left] at (0,0) {$O$};

% 椭圆
\draw[black, thick] (0,0) ellipse [x radius=\a, y radius=\b];
\node[font=\footnotesize, below right] at (2,0) {$a$};
\node[font=\footnotesize, below left] at (-2,0) {$-a$};
\node[font=\footnotesize, right] at (0,1) {$b$};

% 扫描点 P(0, 0.5)
\coordinate (P) at (0, 0.5);
\fill[red] (P) circle [radius=2.5pt];
\node[font=\footnotesize, right, red] at (P) {$P(0,\,\frac{1}{2})$（固定点）};

% 一组过 P 的直线（不同参数 k）
% k=0: y=0.5 (水平线)
\draw[blue!30, thin] (-2.2, 0.5) -- (2.2, 0.5);

% k=0.5: y=0.5x+0.5
\draw[blue!50, thin] ({(-1-0.5)/0.5}, -1) -- ({(1-0.5)/0.5}, 1);
% clip range: x in [-2,2]
\draw[blue!50, thin] (-2, {0.5*(-2)+0.5}) -- (2, {0.5*2+0.5});

% k=1.5
\draw[blue!70, thin] (-2, {1.5*(-2)+0.5}) -- (2, {1.5*2+0.5});

% k=-1
\draw[blue!80, thin] (-2, {-1*(-2)+0.5}) -- (2, {-1*2+0.5});

% k→∞ (竖直线)
\draw[blue!40, dashed, thin] (0, -1.5) -- (0, 2.0);

% 标注不同斜率
\node[font=\tiny, blue!30] at (2.3, 0.5) {$k=0$};
\node[font=\tiny, blue!50] at (2.3, {0.5*2+0.5}) {$k=0.5$};
\node[font=\tiny, blue!70] at (1.7, {1.5*1.7+0.5}) {$k=1.5$};
\node[font=\tiny, blue!80] at (-1.8, {-1*(-1.8)+0.5}) {$k=-1$};
\node[font=\tiny, blue!40] at (0.15, 1.8) {\tiny $k\to\infty$};

% 标注交点（示意几个）
% k=0 与椭圆: y=0.5 → x^2/4 + 0.25=1 → x^2=3 → x=±√3≈±1.732
\coordinate (A1) at (1.732, 0.5);
\coordinate (A2) at (-1.732, 0.5);
\fill[red!60] (A1) circle [radius=1.8pt];
\fill[red!60] (A2) circle [radius=1.8pt];

% k=1.5: y=1.5x+0.5 → x^2/4+(1.5x+0.5)^2=1
% x^2/4 + 2.25x^2+1.5x+0.25=1
% 2.5x^2+1.5x-0.75=0
% x=(−1.5±sqrt(2.25+7.5))/5=(−1.5±sqrt(9.75))/5=(−1.5±3.122)/5
% x1=(−1.5+3.122)/5=0.324, y1=1.5*0.324+0.5=0.987
% x2=(−1.5−3.122)/5=-0.924, y2=1.5*(−0.924)+0.5=-0.887
\coordinate (B1) at (0.324, 0.987);
\coordinate (B2) at (-0.924, -0.887);
\fill[red!60] (B1) circle [radius=1.8pt];
\fill[red!60] (B2) circle [radius=1.8pt];

% 参数说明框
\node[draw, rounded corners, font=\footnotesize, fill=yellow!10, align=left]
  at (0, -1.35)
  {含参直线 $l_k$：$y = kx + \tfrac{1}{2}$，$k$ 为参数\\
   $\Delta > 0 \Leftrightarrow$ 与椭圆相交（讨论 $k$ 的范围）};

\end{tikzpicture}
\end{document}
